\section{Methods}

\subsection{Problem definition and information budget}

Let $c$ denote a city group, $n$ a network, $g$ a matched action group,
$s_g$ its complete simulator state and $\mathcal{A}_g$ the feasible green
stages of the focal signal.  The phase executor constructs this set from
the original SUMO \texttt{tlLogic} and enforces minimum green, yellow and
all-red transitions.  The learner requests a feasible stage, which the
executor translates into the corresponding signal sequence.

The external evaluation uses target offline adaptation.  Available
information comprises (i) a frozen source counterfactual bank, (ii) target
network, demand and signal files, (iii) target counterfactual action groups
from an uncalibrated SUMO model, and (iv) target state observed during
deployment.  The simulator parameters are not calibrated against target field
data, but item (iii) provides labelled target transitions before deployment.
The experiment therefore evaluates target offline adaptation rather than
zero-shot transfer.  Passive traffic observations without action outcomes
would provide a weaker information budget than the one evaluated here.

\subsection{Matched counterfactual action groups}
\label{sec:methods-matched-groups}

After a 60-s warm-up, we save the complete SUMO simulation and all
signal-executor states.  Each focal action $a\in\mathcal{A}_g$ starts from
the same exactly restored state and is applied for one 10-s control interval.
The focal signal then follows PhasePressure for five additional intervals.
All other signals follow PhasePressure throughout the branch.  Candidate cost
is the mean number of active or pending vehicles per controlled lane over
60 one-second samples, with lower values preferred:
\begin{equation}
 Y_g(a)=\frac{1}{H}\sum_{\tau=1}^{H}
 \frac{N_{\mathrm{active}}(\tau)+N_{\mathrm{pending}}(\tau)}{L_n},
 \qquad H=60.
 \label{eq:counterfactual-cost}
\end{equation}
Here $\tau\in\{1,\ldots,H\}$ indexes one-second samples,
$N_{\mathrm{active}}(\tau)$ and $N_{\mathrm{pending}}(\tau)$ are the active
and insertion-pending vehicle counts, respectively, and $L_n$ is the number
of controlled lanes in network $n$.  Thus \Cref{eq:counterfactual-cost} uses a
dimensionless sample count as its summation limit.  Including pending vehicles
prevents a phase from being rewarded for delaying insertion.  The branch collector uses
\texttt{libsumo}, checks replay determinism, fails on teleports, and censors an
entire matched group when any branch has a junction collision.  Group-level
censoring preserves each retained group's complete candidate set.

A generalized pressure policy $\pi_r$, selected using source-only closed-loop
cost, supplies one reference action $r_g=\pi_r(s_g)$.  For the frozen external
models, its specification is \texttt{gp\_d0\_o0p02\_s0}.  Let
$Q_{0.75}(\cdot)$ denote the empirical 75th percentile.  The implementation
uses the numerical floor $\eta_g$, within-group scale $\sigma_g$ and normalized
contrast $z_g(a)$ defined jointly by
\begin{equation}
\begin{aligned}
 \eta_g &= \max\!\left\{10^{-8},
     10^{-6}Q_{0.75}\bigl(\{|Y_g(a)|:a\in\mathcal A_g\}\bigr)\right\},\\
 \sigma_g &= \max\!\left\{
      \max_{a\in\mathcal A_g}Y_g(a)-\min_{a\in\mathcal A_g}Y_g(a),
      \eta_g\right\},\\
 z_g(a) &= \frac{Y_g(a)-Y_g(r_g)}{\sigma_g}.
\end{aligned}
\label{eq:normalized-contrast}
\end{equation}
The reference has $z_g(r_g)=0$.  For a selected action $\widehat a_g$,
evaluation uses the normalized action regret in
\begin{equation}
 \operatorname{Reg}_g(\widehat a_g)
 =\frac{Y_g(\widehat a_g)-\min_{a\in\mathcal{A}_g}Y_g(a)}{\sigma_g}.
 \label{eq:action-regret}
\end{equation}
All reported offline regret values use \Cref{eq:action-regret}.

\begin{assumption}[action-common domain error]
\label{ass:action-common-error}
At one restored state, the candidate cost admits the decomposition
\begin{equation}
 Y_g(a)=b_g+q_g(a)+e_g(a),
 \label{eq:additive-nuisance}
\end{equation}
where $b_g$ is a state-specific term shared by every candidate, $q_g(a)$ is
the action-dependent component intended to transfer, and $e_g(a)$ is the
remaining action-dependent error.
\end{assumption}

\begin{proposition}[action-common cancellation and rank invariance]
\label{prop:exact-cancellation}
Under \Cref{ass:action-common-error}, $b_g$ cancels exactly from every
candidate-versus-reference numerator.  Adding any state-specific constant to
all candidate outcomes leaves the action ordering and regret minimizer
unchanged.  It also leaves the numerical values of $z_g(a)$ unchanged whenever
the within-group range, rather than the relative floor $\eta_g$, determines
$\sigma_g$.
\end{proposition}
\begin{proof}
For any $a\in\mathcal A_g$, subtracting the reference outcome in
\Cref{eq:additive-nuisance} gives
$Y_g(a)-Y_g(r_g)=q_g(a)-q_g(r_g)+e_g(a)-e_g(r_g)$, so $b_g$ cancels.  A common
additive constant cancels from every difference and does not change the raw
within-group range.  Because $\sigma_g>0$, any change in its numerical floor
rescales all contrasts by the same positive factor and cannot change their
ordering.  If the range is the active term in \Cref{eq:normalized-contrast},
$\sigma_g$ itself is unchanged, which also preserves every normalized value.
\end{proof}

This proposition motivates the estimand but does not remove action-dependent
simulator error.  Exact state restoration is therefore a protocol requirement,
not merely a variance-reduction device.

\subsection{Parent-restricted rigid anchor}

Let $\theta$ denote the fitted anchor parameters and define
$A_g(a):=A_\theta(s_g,a)$.  The anchor directly regresses $z_g(a)$ for
non-reference candidates and assigns zero to the reference.  Its 29 declared parents comprise
local total queue, vehicle count, speed and occupancy; candidate and reference
green/red queue; downstream queue and occupancy; service and red pressure;
switch and clearance indicators; and current-phase overlap.  City identifiers,
static target context and unrestricted graph embeddings are excluded from the
direct predictor.  The parent set encodes a predefined modelling hypothesis
about local action mechanisms rather than a structure inferred from field
observations.

The estimator is a \texttt{HistGradientBoostingRegressor} with 100 boosting
iterations, at most 15 leaves, minimum leaf size 12 and $\ell_2$ regularization
10.  Each city domain receives equal total weight, each matched action group
receives equal weight within a domain, and negative versus non-negative
candidate contrasts are reweighted so the dominant sign does not determine
the fit.  Targets are normalized independently by $\sigma_g$.  The action-level
anchor uncertainty $u_{A,g}(a)$ is calibrated from the 90th-percentile absolute
residual on source leave-one-domain predictions.

\subsection{Antisymmetric all-action correction}
\label{sec:methods-pairwise}

The correction model uses every unordered pair $\{a_i,a_j\}\subset
\mathcal{A}_g$, with $i\ne j$ indexing candidate actions.  Let $x_g$ contain
23 action-common state parents and $\phi_g(a)$ contain 40 phase-local action
parents.  With square brackets denoting vector concatenation, the training
design $d_{gij}$ and normalized pair target $t_{gij}$ are
\begin{equation}
 d_{gij}=\left[x_g,\,\phi_g(a_i)-\phi_g(a_j)\right],\qquad
 t_{gij}=\frac{Y_g(a_i)-Y_g(a_j)}{\sigma_g}.
 \label{eq:pair-design}
\end{equation}
Every row defined in \Cref{eq:pair-design} is augmented by its reverse
$(d_{gji},-t_{gij})$.  Let $\psi$ denote the fitted pairwise-model parameters.
At prediction, the regressor $f_\psi$ is explicitly antisymmetrized,
\begin{equation}
 \widehat\delta_{gij}=\tfrac{1}{2}\left[
 f_\psi(x_g,\phi_g(a_i)-\phi_g(a_j))
 -f_\psi(x_g,\phi_g(a_j)-\phi_g(a_i))\right],
 \label{eq:antisymmetry}
\end{equation}
and converted first to the unshifted all-pair score $\widetilde P_g$ and then
to the reference-shifted score $P_g(a):=P_\psi(s_g,a)$:
\begin{equation}
\begin{aligned}
 \widetilde P_g(a_i)&=\frac{1}{K_g-1}
       \sum_{j\neq i}\widehat\delta_{gij},
 &K_g&=|\mathcal{A}_g|,\\
 P_g(a_i)&=\widetilde P_g(a_i)-\widetilde P_g(r_g).
\end{aligned}
 \label{eq:borda-score}
\end{equation}
Thus $P_g(r_g)=0$.  Pair-cycle inconsistency augments the 90th-percentile
residual in the action-level pairwise uncertainty $u_{P,g}(a)$.

\begin{proposition}[antisymmetry and ranking recovery]
\label{prop:pairwise-recovery}
\Cref{eq:antisymmetry} gives
$\widehat\delta_{gji}=-\widehat\delta_{gij}$.  If
$\widehat\delta_{gij}=z_g(a_i)-z_g(a_j)$ for every pair, then
\begin{equation}
\begin{aligned}
 \widetilde P_g(a_i)&=\frac{K_g}{K_g-1}
 \left[z_g(a_i)-\overline z_g\right],\\
 P_g(a_i)&=\frac{K_g}{K_g-1}
 \left[z_g(a_i)-z_g(r_g)\right].
\end{aligned}
 \label{eq:pairwise-recovery}
\end{equation}
Here $\overline z_g=K_g^{-1}\sum_{a\in\mathcal A_g}z_g(a)$.  The positive
scale preserves $\arg\min_a z_g(a)$.
\end{proposition}
\begin{proof}
Swapping $i$ and $j$ in \Cref{eq:antisymmetry} negates the expression.  Under
the stated exact-pair condition, summing in \Cref{eq:borda-score} yields
\Cref{eq:pairwise-recovery}; subtracting a common mean and multiplying by a
positive constant cannot change the minimizer.
\end{proof}

The pairwise model uses the same gradient-boosting family as the anchor.  Its
state parents comprise local aggregate state, time, queue concentration and
signal-graph neighbourhood summaries; its action parents include phase
duration, queue, occupancy, downstream load, service pressure, transition and
graph-action summaries.  Context values are used only to calibrate support and
uncertainty and are excluded from the direct prediction features.

\subsection{Anchored interpolation and local cap}
\label{sec:methods-blend}

For one action group, the deployed score is
\begin{equation}
 S_g(a)=(1-\alpha_g)A_g(a)+\alpha_gP_g(a),
 \qquad 0\leq\alpha_g\leq1.
 \label{eq:anchored-blend}
\end{equation}
Candidate selectors include constant choices
$\alpha_g\in\{0,0.1,\ldots,1\}$ and local confidence-capped weights.  For the
local form, for any action-indexed vector $v_g$ define
$\operatorname{range}(v_g)=\max_{a\in\mathcal A_g}v_g(a)-
\min_{a\in\mathcal A_g}v_g(a)$.  Let
\begin{equation}
\begin{aligned}
 a_{A,g}&=\arg\min_{a\in\mathcal A_g}A_g(a),
 &a_{P,g}&=\arg\min_{a\in\mathcal A_g}P_g(a),\\
 h_g &= \max\{\operatorname{range}(A_g),
                  \operatorname{median}_{a}u_{A,g}(a),10^{-6}\},\\
 D_g &= \operatorname{range}(P_g-A_g)/h_g,\\
 M_g &= \max\{P_g(a_{A,g})-P_g(a_{P,g}),0\}/h_g,\\
 U_g &= \{u_{P,g}(a_{A,g})+u_{P,g}(a_{P,g})\}/h_g.
\end{aligned}
\label{eq:local-diagnostics}
\end{equation}
Here $h_g$ is the local score scale, $D_g$ the normalized score disagreement,
$M_g$ the pairwise preference margin and $U_g$ its uncertainty sum.  With
perturbation cap $\kappa$ and confidence multiplier $\zeta$, the diagnostics
in \Cref{eq:local-diagnostics} determine the local weight
\begin{equation}
 \alpha_g=\min\{1,\kappa/\max(D_g,10^{-12})\}
           \mathbf{1}\{M_g\geq\zeta U_g\},
 \label{eq:local-alpha}
\end{equation}
where $\mathbf{1}\{\cdot\}$ is the indicator function.

\begin{proposition}[blend and decision-error bounds]
\label{prop:blend-bounds}
For a group-constant $\alpha_g$, if the non-negative constants
$\varepsilon_A$ and $\varepsilon_P$ bound the anchor and pairwise uniform
score errors, respectively, so that
$\|A_g-z_g\|_\infty\leq\varepsilon_A$ and
$\|P_g-z_g\|_\infty\leq\varepsilon_P$, then
\begin{equation}
 \|S_g-z_g\|_\infty
 \leq(1-\alpha_g)\varepsilon_A+\alpha_g\varepsilon_P.
 \label{eq:blend-error-bound}
\end{equation}
For $\widehat a_g=\arg\min_{a\in\mathcal A_g} S_g(a)$ and
$a_g^*=\arg\min_{a\in\mathcal A_g} z_g(a)$,
\begin{equation}
 z_g(\widehat a_g)-z_g(a_g^*)
 \leq2\|S_g-z_g\|_\infty.
 \label{eq:decision-bound}
\end{equation}
Here $\|v_g\|_\infty=\max_{a\in\mathcal A_g}|v_g(a)|$.
\end{proposition}
\begin{proof}
The error bound in \Cref{eq:blend-error-bound} follows from convexity of the
sup norm.  The decision bound in \Cref{eq:decision-bound} follows
by adding and subtracting $S_g(\widehat a_g)$ and $S_g(a_g^*)$, using the
minimizing property of $\widehat a_g$, and bounding each remaining discrepancy
by $\|S_g-z_g\|_\infty$.
\end{proof}
These are deterministic conditional bounds, not guarantees that either fitted
error is small in a new city.

\begin{corollary}[bounded local perturbation]
\label{cor:bounded-local-perturbation}
Whenever the confidence indicator in \Cref{eq:local-alpha} passes,
\begin{equation}
 \operatorname{range}(S_g-A_g)
 =\alpha_g\operatorname{range}(P_g-A_g)\leq \kappa h_g.
 \label{eq:cap-bound}
\end{equation}
When it fails, $S_g=A_g$.  Thus \Cref{eq:cap-bound} bounds score
perturbation relative to the anchor, but it is not a queue-stability or safety
guarantee.
\end{corollary}

\subsection{Seed-blocked target selection and full refit}

For each external city, we retain all valid groups from adaptation seeds
5057 and 6067.  In fold one, both models are fitted on the source bank and
target seed 5057; target seed 6067 evaluates every candidate $\alpha_g$.
Fold two reverses the target seeds.  An admissible candidate must match or
improve the anchor's mean OOF regret and regress by at most 0.1 in either
complete-seed fold.  The selector minimizes mean OOF regret among admissible
candidates, with deterministic tie breaking.  Los Angeles selects
$(\kappa,\zeta)=(1,0.25)$ and Jinan selects constant $\alpha_g=0.9$.

After the choice is frozen, the model is fitted once on the source bank plus
all target adaptation groups: 165 groups for Los Angeles and 911 for Jinan.
The OOF models are discarded.  CFCMT and the learned baselines share the
same target group identifiers and full-refit call.  Serialized models are
reloaded and checked before deployment.  The seed-8171 evaluation cache is required
to be absent at the joint freeze and the diagnostic-only seed 7079 is
permanently excluded from V43 selection and confirmation.

\subsection{Closed-loop controller and comparators}

At each 10-s decision time, the runtime enumerates every currently feasible
phase, evaluates its frozen 60-s action score and requests the phase with the
minimum deployed score $S_g(a)$.  The phase executor enforces physical
transition timing.
No uncertainty guard, pressure fallback or target closed-loop tuning is active
in the final V43 policy.  The pressure reference is recomputed only to create
the zero coordinate required by the contrast model.

The same-protocol comparators are fixed time, MaxPressure, PhasePressure,
simulator-only contrast MPC, the rigid anchor alone and a dense residual.  The
dense residual receives the same matched action contrasts and refit budget
but allows unrestricted feature/context mixing in a dense residual world
model.  It is labelled H2O$^+$-style because it combines a simulator with a
learned residual.  It does not reproduce the original H2O$^+$ algorithm,
optimizer or online training loop.  Because both methods use matched action
contrasts, subtracting the reference action alone cannot explain any CFCMT
advantage.

Official RESCO and LibSignal controllers retain their upstream state, reward,
timing and training definitions.  These runs assess execution under the
native protocols rather than the paired CFCMT comparison.

\subsection{Legacy near-target-only source-contribution test}

The original source-contribution test uses the pre-existing
\texttt{causal\_target\_only} model family.  This comparator has the same causal
action-advantage features, target domain and target adaptation groups as the
corresponding baseline artifact.  It sets prior group strength to zero and the
maximum target weight to one.  Both external artifacts had more than four
target groups and selected target weights numerically equal to one.  A source
submodel is still fitted and retained; its coefficient in the deployed mean,
uncertainty and trust predictions is zero, but its prior can resolve exact
action ties.  We therefore designate this frozen arm the legacy near-target-only
comparator.  Its retained tie-breaking prior distinguishes it from a strict
zero-source-row estimator.

The first deployment of this frozen comparator, V90, reused the eight V43
scenario--seed cells and was specified only after the original seven-policy
outcomes were known.  V90 is therefore post-hoc.  The subsequent V91 protocol
was frozen before any V91 outcome, with neither model refitted.  A deterministic
generator produced 64 seeds disjoint from every seed-valued field in V43 and
V89 and from the two V90 closed-loop seeds.  Four scenarios and two policies
gave 512 required rollout identities, all of which completed.

For policy $\pi$, city $c$, demand scenario $m$ in the city-specific set
$\mathcal M_c$ and seed index $k$, let $W^\pi_{c,m,k}$ be the all-departed
mean waiting time.  With the external city set
$\mathcal C=\{\text{Los Angeles},\text{Jinan}\}$, the scenario and equal-city
means are
\begin{equation}
 \bar W^{\pi}_{c,k}=|\mathcal M_c|^{-1}
   \sum_{m\in\mathcal M_c} W^{\pi}_{c,m,k},\qquad
 \bar W^{\pi}_{k}=|\mathcal C|^{-1}
   \sum_{c\in\mathcal C} \bar W^{\pi}_{c,k}.
 \label{eq:waiting-aggregation}
\end{equation}
Using the aggregation in \Cref{eq:waiting-aggregation}, the paired relative
waiting effect is
\begin{equation}
 \rho_k=\frac{\bar W^{\mathrm{CFCMT}}_k-
              \bar W^{\mathrm{legacy}}_k}
             {\bar W^{\mathrm{legacy}}_k},
 \label{eq:source-contribution}
\end{equation}
where $\pi=\mathrm{legacy}$ denotes the frozen legacy near-target-only
comparator and negative values favour CFCMT.  The primary estimate is the mean of
$\rho_k$ over 64 seeds, with 20,000 paired-seed bootstrap replicates.  The
predeclared joint gate required a non-positive mean and interval upper bound,
non-positive means in both cities, at most 5\% worst-seed regression, at least
50\% improved seeds, one-sided exact sign-test probability value
$p_{\mathrm{sign}}\leq0.05$, no excess
teleports or collision incidents, all identities and no failed-seed exclusion.
Failure of any component rejects the source-contribution confirmation.

\subsection{Post-V91 target-offline conditional source selection}

V98 is a separate post-V91 successor and is not part of the frozen V43 method.
Target models were refitted after repairing topology context, with one candidate
component for each of seven source-city groups.  Five scenario-and-seed-blocked
folds evaluated source masses $\{0,0.25,0.5,0.75,1\}$ on a balanced B100 subset
of target counterfactual groups.  A source component was admitted only when its
mean normalized-regret improvement exceeded 0.5\%, its one-sided 95\% upper
confidence bound was below zero and its worst fold regressed by less than 1\%.
Failure of any gate selected the exact target-only model.  This rule returned
source mass zero for Los Angeles and Hangzhou source mass one for Jinan before
the fresh confirmation rollouts were observed.

The strict comparator uses the corrected
\texttt{causal\_target\_only\_v2} estimator and fits target rows only.  The fresh
Jinan confirmation used 56 deterministic seeds produced by NumPy's 64-bit
Permuted Congruential Generator (PCG64), excluded from V93 and the earlier V98
confirmation.  Three
demand scenarios and two policies gave 336
completed rollouts.  Waiting was averaged equally across the three scenarios
within each seed; uncertainty used a paired-seed bootstrap, and the paired
Wilcoxon statistic was computed over the 56 seed means.

The selector uses no target closed-loop outcome, although the full V98
protocol incorporates earlier closed-loop development.  Both Jinan arms share
a pressure-guard risk multiplier of 0.5 inherited from that development.
The strict zero-source-row comparator also has lower model capacity than the
source-aware estimator.  V98 therefore evaluates the frozen conditional
protocol without isolating source-row contribution at fixed architecture.

\subsection{Separate post-freeze execution-layer stress test}

The V86--V89 track is a distinct successor experiment on
\texttt{jinan\_3x4\_real}.  Its originator predicts a 120-s target-simulator cost
from target-state and action features; the executor permits at most one
network-wide override and requires at least 120 s between effective phase
switches.  The method was selected only from 54 disclosed adaptive seeds.
Unlimited cooldown stays and a one-stay-per-cooldown variant were evaluated on
the same adaptive partition and could not replace the legacy executor.

Before any V89 rollout, a deterministic generator selected 64 new seeds after
excluding the 54 adaptive seeds and eight separately sealed V85 seeds.  Each
seed paired the frozen executor with PhasePressure under the same target
scenario.  The preregistered joint rule required non-positive mean relative
waiting difference, a non-positive paired-bootstrap 95\% upper bound, no more
than 5\% worst-seed regression, at least 50\% improved seeds, one-sided exact
sign-test $p_{\mathrm{sign}}\leq0.05$, at least 75\% active seeds, no aggregate excess in
teleports or collision incidents, closed phase accounting and all 128 rollout
identities.  Twenty thousand paired bootstrap replicates were fixed in the
protocol.  No failed seed was excluded, and V85 remained sealed.  This track is
reported as a stress test and is not part of V43 method selection.

\subsection{Networks, splits and experimental execution}

Development uses 18 networks in seven groups.  Source model seeds are 2027 and
3037 and the excluded development evaluation seed is 4047.  The successful
development stage uses up to 100 target-style groups inside each held-out
source-city exercise.  External confirmation uses one Los Angeles network and
one Jinan geometry with three demand files.  Jinan scenarios are weighted
equally before city aggregation, preventing that city from receiving triple
the weight of Los Angeles.

Closed-loop evaluation runs four scenarios, two seeds and seven policies for
56 paired rollouts.  Every rollout uses SUMO 1.22.0, \texttt{libsumo}, a 60-s
warm-up and a 3,600-s measurement horizon.  The experiment was parallelized
across cluster workers, while numerical libraries were restricted to one
thread per process to avoid nested oversubscription.  Each rollout records its
result and \texttt{tripinfo} hashes, source snapshot identity, model identity,
SUMO version and metric population.

The fresh source-contribution confirmation uses the same four scenarios,
simulator version, warm-up, horizon, action executor and metric implementation.
Only the two frozen policies and 64 fresh seeds differ, producing 512 rollouts.
Six process-isolated shards ran on six nodes with 22 workers per node.  Full
rollout records remain on the execution host; six complete shard summaries and
the launch and audit manifests are materialized with the manuscript evidence.

\subsection{Outcomes, aggregation and uncertainty}

The preregistered closed-loop primary is mean vehicle waiting time.  The
implementation computes it over all departed vehicles, including unfinished
\texttt{tripinfo} records whose waiting and duration are accrued to the fixed
horizon.  We additionally report completed-only waiting, system vehicle-hours,
unfinished fraction, completion ratio, throughput, mean queue per lane,
teleports, emergency stops and junction collisions.  Together, these outcomes
account for withheld departures and unfinished trips that could otherwise
make a controller appear effective.

Scenario-seed results are first averaged within scenario, scenarios are
averaged equally within city and the two external cities are averaged equally.
All policy contrasts remain paired on city, scenario and seed.  Offline
uncertainty uses 10,000 paired action-group bootstrap replicates.  Closed-loop
uncertainty uses 10,000 paired hierarchical replicates that sample city, then
scenario, then seed.  With only two independent target cities, these intervals
describe sensitivity of the observed matrix; they are not estimates of a
general population of cities.  No multiplicity-adjusted claim or universal
city-level significance is made.

The V91 source-contribution interval instead resamples the 64 paired seed
statistics in \Cref{eq:source-contribution} 20,000 times.  The exact sign test
uses the number of seeds with $\rho_k<0$.  City-specific relative differences
are secondary diagnostics because the independent city count remains two.

Raw SUMO collision events are consolidated by participant, lane and time into
unique junction incidents.  Teleports are a hard validity failure; collisions
are retained as outcomes in closed loop and compared pairwise.  A valid matrix
requires all 56 authorized records, exact provenance identities, complete
population accounting, readable \texttt{tripinfo} hashes and zero teleports.
The V91 integrity audit analogously requires all 512 frozen identities and
retains every seed regardless of performance or incident status.

\subsection{Meaning of ``causal''}

CFCMT follows two causal design commitments.  Candidate actions intervene
from a common restored simulator state, and direct predictors are restricted
to declared local parent sets.  \Cref{prop:exact-cancellation} is algebraic under
\Cref{ass:action-common-error}, while \Cref{prop:pairwise-recovery,prop:blend-bounds}
and \Cref{cor:bounded-local-perturbation} concern score construction.  The
study does not identify a causal graph from observational traffic, prove that
the declared parents are invariant in the field, or estimate effects of signal
interventions in Los Angeles or Jinan.  ``Causal'' therefore describes the
model and counterfactual simulation design, not field causal identification.
